% TOP_PROBLEM: {"id": 10000015, "title": "Roughly transitive graphs versus homogeneous spaces", "category_id": 3, "category": {"id": 3, "name": "graph_theory", "display_name": "Graph Theory", "description": "Problems involving graphs, networks, and their properties.", "slug": "graph-theory", "order_index": 3, "created_at": "2026-07-31T15:26:25.671Z"}, "status": "open", "problem_number": "AMR-099-0015", "set_id": 13}
% FIRST_POSTED: 2026-10-01
% ENTRY_KIND: statement_only
% UnsolvedMath ID 10000015; source catalogue code AMR-099-0015
% !TeX program = lualatex
% Definition and sources only; this file is not an LLM solution attempt.
\documentclass[11pt,a4paper]{article}
\usepackage[margin=25mm]{geometry}
\usepackage{fontspec}
\setmainfont{lmroman10-regular.otf}[BoldFont=lmroman10-bold.otf,
 ItalicFont=lmroman10-italic.otf,BoldItalicFont=lmroman10-bolditalic.otf]
\usepackage{amsmath,amssymb,mathtools}
\usepackage{unicode-math}
\setmathfont{latinmodern-math.otf}
\AtBeginDocument{\let\setminus\smallsetminus}
\usepackage{xurl,hyperref}
\hypersetup{colorlinks=true,urlcolor=blue}
\setcounter{secnumdepth}{0}
\setlength{\parindent}{0pt}
\setlength{\parskip}{5pt}
\setlength{\emergencystretch}{3em}
\begin{document}
\section*{Roughly transitive graphs versus homogeneous spaces}\label{10000015}
{\small UnsolvedMath ID 10000015 · AMR-099-0015 · Graph Theory · AMR Open Problem Lists}
\subsection{Definitions and mathematical statement}
For metric spaces $X,Y$, a map $f:X\to Y$ is an $(A,B)$-quasi-isometry, with $A\ge1$ and $B\ge0$, if
\[
A^{-1}d_X(x,x')-B\le d_Y(f(x),f(x'))\le A d_X(x,x')+B
\]
for all $x,x'\in X$, and every point of $Y$ is within distance $B$ of $f(X)$. The term rough isometry in this problem permits these multiplicative and additive constants.

Call a connected graph $G$, with its graph metric, uniformly roughly transitive if some fixed finite $A,B$ have the following property: for every pair of vertices $x,y$, there is an $(A,B)$-quasi-isometry $f_{x,y}:G\to G$ with $f_{x,y}(x)=y$. A metric space $M$ is homogeneous if its group of distance-preserving bijections acts transitively on its points.

Does there exist an infinite uniformly roughly transitive graph which is not quasi-isometric to any homogeneous metric space? Equivalently, is every such graph quasi-isometric to some homogeneous space, with no requirement that this space itself be a graph or a Cayley graph?

Uniformity of $A,B$ over the choice of $x,y$ is part of the hypothesis. Allowing constants to depend arbitrarily on that pair would give a different and much weaker notion. A proposed homogeneous comparison space must admit actual isometries between all pairs of points, while the original graph only supplies approximate self-symmetries. The question asks whether the approximate symmetry can always be replaced, at bounded metric distortion, by exact homogeneity.
\subsection{Short English statement}
Is every graph with uniform approximate transitivity coarsely equivalent to a space with exact transitive isometries?
\subsection{Sources}
\begin{itemize}
\item[S1] ULAM.AI and the UnsolvedMath Contributors, supplied problems.json, record 10000015 (AMR-099-0015), AMR Open Problem Lists. Catalogue identity and source transcription; CC BY 4.0.
\url{https://huggingface.co/datasets/ulamai/UnsolvedMath}.
\item[S2] Benjamini - Coarse Geometry and Randomness (Saint-Flour notes), Open problem 2.3, PDF page 22. Original source indexed by AMR-099-0015.
\url{https://www.wisdom.weizmann.ac.il/~itai/stflouraug24.pdf#page=22}.
\item[S3] I. Benjamini, Coarse Geometry and Randomness, primary Saint-Flour notes; the numbered question and its surrounding definitions.
\url{https://arquivo.pt/noFrame/replay/20201231041548id_/http://www.wisdom.weizmann.ac.il/~itai/stflouraug24.pdf}.
\item[S4] I. Benjamini, Euclidean vs. Graph Metric, primary numbered question and conventions.
\url{https://arquivo.pt/noFrame/replay/20201231041538id_/http://www.wisdom.weizmann.ac.il/~itai/erd100.pdf}.
\end{itemize}
\end{document}
